\BourbakiExerciseGroup

\begin{enumerate}
\def\labelenumi{\arabic{enumi})}
\tightlist
\item
  设 K 是特征为 p \textgreater{} 0 的域，E 是 K 的一个扩张；E 是一个高度 \(\leqslant 1\) 的、\(\mathrm{K}(\mathrm{E}^{p})\) 的 p-根式扩张。E 在 \(\mathrm{K}(\mathrm{E}^{p})\) 上的 p-基的基数（参见 V，第 103 页，定理 2）称为 E 在 K 上的不完全度。若 \(\mathrm{K}(\mathrm{E}^{p}) = \mathrm{E}\) ，则称 E 在 K 上相对完全；若 E 是完全域，则它在其每个子域上都相对完全。K 的（绝对）不完全度是 K 在其素子域 \(F_{p}\) 上的不完全度，亦即 K 的一个（绝对）p-基的基数。
\end{enumerate}

\begin{enumerate}
\def\labelenumi{\alph{enumi})}
\item
  若 B 是 E 的一个 \(p\) -基（在 \(\mathbf{K}(\mathbf{E}^p)\) 上），证明对每个整数 \(k > 0\) ，我们有 \(\mathbf{E} = \mathbf{K}(\mathbf{E}^{p^k})(\mathbf{B})\) 。
\item
  假设 \(\mathbf{E} \subset \mathbf{K}^{p^{-n}}\) 对某个整数 \(n > 0\) 成立。证明：要使次数 \([\mathbf{E}:\mathbf{K}]\) 有限，必须而且只需不完全度 \(m_0\) —— \(\mathbf{E}\) 在 \(\mathbf{K}\) 上的 —— 为有限；此时 \(m_0\) 是集合 \(\mathbf{S}\) （满足 \(\mathbf{E} = \mathbf{K}(\mathbf{S})\) ）的最小基数。证明：若 \(m_k\) 是 \(\mathbf{K}(\mathrm{E}^{p^k})\) 在 \(\mathbf{K}\) 上的不完全度，则 \(m_{k + 1} \leqslant m_k\) 对所有 \(k\) 成立，且若 \(f = \sum_{k} m_k\) ，则 \([\mathbf{E}:\mathbf{K}] = p^f\) 。
\end{enumerate}

\begin{enumerate}
\def\labelenumi{\arabic{enumi})}
\setcounter{enumi}{1}
\tightlist
\item
  \begin{enumerate}
  \def\labelenumii{\alph{enumii})}
  \tightlist
  \item
    设 E 是特征为 p \textgreater{} 0 的域 K 的扩张，F 是 E 的扩张。证明：若 B 是 E 在 \(\mathrm{K}(\mathrm{E}^{p})\) 上的 p-基，C 是 F 在 \(\mathrm{E}(\mathrm{F}^{p})\) 上的 p-基，则存在 F 在 \(\mathrm{K}(\mathrm{F}^{p})\) 上的一个包含于 \(B \cup C\) 的 p-基。
  \end{enumerate}
\end{enumerate}

\begin{enumerate}
\def\labelenumi{\alph{enumi})}
\setcounter{enumi}{1}
\tightlist
\item
  假设 B 有限，且 F 是 K 的有限次代数扩张。证明 \([E:K(E^{p})]\geqslant[F:K(F^{p})]\) 。若进一步 F 在 E 上可分，证明 \([E:K(E^{p})]=[F:K(F^{p})]\) 。
\end{enumerate}

\begin{enumerate}
\def\labelenumi{\arabic{enumi})}
\setcounter{enumi}{2}
\item
  设 K 是一个非完全域，E 是 K 的有限次代数扩张。要使存在元素 \(x \in E\) 使得 \(\mathrm{E} = \mathrm{K}(x)\) ，必须而且只需 E 在 K 上的不完全度（习题 1）为 0 或 1。（要看出该条件是充分的，注意：若 \(E_{s}\) 是 K 在 E 中的相对可分闭包，我们有 \(\mathrm{E} = \mathrm{E}_{s}(\alpha)\) 且 \(\mathrm{E}_{s} = \mathrm{K}(\beta)\) ；若 \(\alpha \notin E_{s}\) ，考虑 \(\alpha + c\beta\) 在 K 上的共轭元，其中 \(c \in K\) 。）
\item
  设 K 是特征为 p \textgreater{} 0 的域，E 是 K 的有限次代数扩张。若 r \textgreater{} 0 是 E 在 K 上的不完全度（习题 1），证明 r 是使 \(\mathrm{E} = \mathrm{K}(\mathrm{S})\) 的集合 S 的最小基数。（要看出 E 可由它的 r 个元素生成，注意：若 \(E_{s}\) 是 K 在 E 中的相对可分闭包，则存在 E 的 r 个元素 \(a_{i}\) （ \(1 \leqslant i \leqslant r\) ）使得 \(\mathrm{E} = \mathrm{E}_{s}(a_{1}, \ldots, a_{r})\) ，并利用习题 3。）
\item
  设 K 是一个域，E 是 K 的有限次代数扩张。
\end{enumerate}

\begin{enumerate}
\def\labelenumi{\alph{enumi})}
\item
  假设 K 的特征为 p \textgreater{} 0。证明：若 E 在 K 上的不完全度（习题 1）\textgreater{} 1，则存在无穷多个不同的域 F，使得 \(K \subset F \subset E\) （归结为 \(\mathrm{K}(\mathrm{E}^{p}) = \mathrm{K}\) 的情形，并考虑域 \(\mathrm{K}(a + \lambda b)\) ，其中集合 \(\{a, b\}\) 在 K 上是 p-自由的，且 \(\lambda \in K\) ）。
\item
  由此得出结论：要使满足 \(K \subset F \subset E\) 的域 F 只有有限多个，必须而且只需 E 在 K 上的不完全度为 \(\leqslant 1\) 。
\end{enumerate}

\begin{enumerate}
\def\labelenumi{\arabic{enumi})}
\setcounter{enumi}{5}
\tightlist
\item
  设 K 是特征为 p \textgreater{} 0 的域，E 是 K 的有限次代数扩张，N 是由 E 生成的 K 的拟伽罗瓦扩张。
\end{enumerate}

\begin{enumerate}
\def\labelenumi{\alph{enumi})}
\item
  设 \(E_{r}\) 与 \(E_{s}\) （相应地 \(N_{r}\) 与 \(N_{s}\) ）是包含在 E（相应地 N）中的 K 的最大 p-根式扩张与最大可分扩张。证明：要使 \(\mathrm{E} = \mathrm{E}_{r}(\mathrm{E}_{s})\) （此时 E 同构于 \(E_{s} \otimes_{K} E_{r}\) ），必须而且只需 \(E_{r} = N_{r}\) 。
\item
  若 E 在 K 上的不完全度为 1，则 \(\mathrm{E} = \mathrm{E}_{r}(\mathrm{E}_{s})\) 的充分必要条件是 N 在 K 上的不完全度为 1。由此推出：若 N 是 K 的拟伽罗瓦扩张，且其在 K 上的不完全度等于 1，则对 N 的每个子扩张 E 都有 \(\mathrm{E} = \mathrm{E}_{r}(\mathrm{E}_{s})\) （参见习题 2，b））。
\item
  反之，若 N 是 K 的有限次拟伽罗瓦扩张，且其在 K 上的不完全度 \textgreater1，证明：若 \(N_{r} \neq N\) ，则存在 \(t \in N\) 使得对 K 的扩张 \(\mathrm{E} = \mathrm{K}(t)\) 我们有 \(\mathrm{E} \neq \mathrm{E}_{r}(\mathrm{E}_{s})\) （若 a、b 是 N 在 \(\mathrm{K}(\mathrm{N}^{p})\) 上的一个 p-基中的两个元素，且 \(x \in N\) 在 K 上可分并满足 \(x \notin K\) ，取 \(t = a + bx\) ）。
\end{enumerate}

